Hiện nay, các bộ ngữ liệu đối sánh chuẩn (benchmark) cho việc phân tích sắc thái cảm xúc tiếng Việt khá khiêm tốn về quy mô: ViLexNorm chỉ có hơn 10 ngàn cặp (câu - nhãn), UIT-VSMEC khoảng 7 ngàn cặp; ngay cả bộ dữ liệu ViGoEmotions mới nhất vốn được xem là lớn nhất hiện có, cũng chỉ dừng ở khoảng 21 ngàn cặp trải trên 28 lớp cảm xúc, được gán nhãn theo hướng đa nhãn.

So sánh với các kho ngữ liệu quốc tế như YouNiverse \parencite{ribeiro2021youniverse} (hơn 8,6 tỷ lượt tương tác bình luận) hay Douban-Dushu \parencite{zhao2018lsicc} (hơn 37 triệu mẫu đánh giá từ nền tảng đọc sách), quy mô dữ liệu tiếng Việt hiện có bộc lộ khoảng cách rất lớn. Trong khi đó, để phát huy tối đa năng lực biểu diễn ngữ nghĩa của các kiến trúc học sâu tiên tiến như Transformer, tập dữ liệu đầu vào cần thỏa mãn yêu cầu về quy mô lẫn chất lượng. Tính "lớn" ở đây không chỉ nằm ở số lượng, mà còn thể hiện qua sự đa dạng trong cấu trúc từ vựng, khả năng bao quát đa miền (multi-domain) và độ phong phú của các tầng ngữ nghĩa biểu đạt.

Xuất phát từ thực trạng này, mục tiêu cốt lõi của giai đoạn thu thập dữ liệu là thu về hàng triệu dòng bình luận thô từ YouTube, xem đây là tài liệu chuẩn để hình thành một tập dữ liệu chuẩn, vừa phong phú trên nhiều lĩnh vực, vừa phản ánh trung thực sự phức tạp và các đặc trưng nhiễu vốn có của văn phong trực tuyến tại Việt Nam. Qua đó, nghiên cứu kỳ vọng kiến tạo được một không gian dữ liệu đủ sâu rộng, đáp ứng trọn vẹn yêu cầu huấn luyện các mạng nơ-ron đánh giá cảm xúc ở cấp độ vi mô với độ chuẩn xác và tính linh hoạt vượt trội.

Để hiện thực hóa mục tiêu này, toàn bộ quy trình được triển khai qua hai bước chính.

\textbf{Bước 1}, Danh sách hạt giống gồm các kênh phổ biến tại Việt Nam trải rộng trên mọi hạng mục mà nền tảng này phân loại, dựa vào sự hỗ trợ của hệ thống thống kê trung gian uy tín \textit{socialblade} \parencite{socialblade}, sau đó áp dụng thuật toán \ref{alg:seed_channel_expansion} để mở rộng tập danh sách hạt giống này.

\begin{algorithm}[H]
\caption{Thuật toán mở rộng tập dữ liệu kênh hạt giống dựa trên đề xuất của YouTube}
\label{alg:seed_channel_expansion}
\begin{algorithmic}[1]
\Require $S_{channel}$: Tập hạt giống ban đầu (danh sách các kênh)
\Ensure $S_{channel}$: Tập danh sách kênh sau khi đã được mở rộng toàn diện

\State $Q \gets S_{channel}$ \Comment{Khởi tạo hàng đợi $Q$ chứa các kênh cần duyệt}

\While{$Q \neq \emptyset$} \Comment{Lặp cho tới khi duyệt hết danh sách kênh (hàng đợi rỗng)}
    \State $i \gets \text{Lấy (Pop) một kênh từ đầu hàng đợi } Q$
    \State $V_i \gets \text{Lấy danh sách các video thuộc kênh } i$

    \For{\textbf{each} video $j \in V_i$} \Comment{Duyệt từng video $j$ của kênh $i$}
        \State $C_{rec} \gets \text{Lấy 3 kênh đề xuất mới từ YouTube liên quan đến video } j$

        \For{\textbf{each} kênh $c \in C_{rec}$}
            \If{$c \notin S_{channel}$} \Comment{Đảm bảo kênh đề xuất chưa từng có trong tập hạt giống}
                \State $S_{channel} \gets S_{channel} \cup \{c\}$ \Comment{Bổ sung ngay vào tập hạt giống}
                \State $\text{Đẩy (Push) kênh } c \text{ vào cuối hàng đợi } Q$ \Comment{Để tiếp tục duyệt video của kênh này ở các lượt sau}
            \EndIf
        \EndFor
    \EndFor
\EndWhile

\State \Return $S_{channel}$
\end{algorithmic}
\end{algorithm}

\textbf{Bước 2}, duyệt toàn bộ video thuộc các kênh có trong danh sách đường dẫn hạt giống (URLs) để thu thập siêu dữ liệu (metadata) video cùng hệ thống bình luận đi kèm. Nhằm phục vụ việc đánh giá sắc thái cảm xúc trên đa lĩnh vực và kiểm soát tối ưu nguồn thông tin trực tuyến, dữ liệu thu về tổ chức theo mô hình phân cấp (hierarchical structure), tương tự cách tiếp cận mà \textcite{ding2026scope} đề xuất giúp tách bạch rõ ràng giữa thực thể video (nút cha) và các luồng thảo luận (nút con), đồng thời bảo toàn cấu trúc phản hồi nhiều cấp (nested replies) của người dùng.

Các trường thông tin cùng cấu trúc lưu trữ chi tiết của hệ thống dữ liệu này được trình bày cụ thể trong Bảng \ref{tab:video_metadata} và Bảng \ref{tab:comment_metadata}.

\begin{table}[H]
\centering
\caption{Các trường siêu dữ liệu cốt lõi của Video}
\label{tab:video_metadata}
\begin{tabular}{|l|l|p{8cm}|}
\hline
\textbf{Tên trường} & \textbf{Kiểu dữ liệu} & \textbf{Mô tả chi tiết} \\ \hline
\texttt{video\_id}   & Varchar (PK)          & Mã định danh duy nhất của video do YouTube cấp. \\ \hline
\texttt{title}      & Text                  & Tiêu đề chính thức của video. \\ \hline
\texttt{category}   & Varchar               & Tên danh mục nội dung (ví dụ: edu, music, news, \dots). \\ \hline
\texttt{channel}    & Varchar               & Tên của kênh YouTube đăng tải video. \\ \hline
\texttt{video\_type}    & Varchar               & Loại video (gồm: static, short, stream). \\ \hline
\end{tabular}
\end{table}

\begin{table}[H]
\centering
\caption{Cấu trúc phân cấp các trường thông tin của Bình luận}
\label{tab:comment_metadata}
\begin{tabular}{|l|l|p{8cm}|}
\hline
\textbf{Tên trường} & \textbf{Kiểu dữ liệu} & \textbf{Mô tả chi tiết} \\ \hline
\texttt{comment\_id} & Varchar (PK)          & Mã định danh duy nhất của bình luận. \\ \hline
\texttt{video\_id}   & Varchar (FK)          & Mã định danh của video chứa bình luận (khóa ngoại liên kết với Bảng \ref{tab:video_metadata}). \\ \hline
\texttt{parent\_id}  & Varchar (FK)          & Mã định danh của bình luận gốc (chỉ xuất hiện đối với các phản hồi, nhận giá trị \texttt{NULL} nếu là bình luận cấp 1). \\ \hline
\texttt{comment}     & Text                  & Nội dung văn bản thô của bình luận do người dùng nhập. \\ \hline
\texttt{created\_at} & Datetime              & Mốc thời gian bình luận được khởi tạo. \\ \hline
\texttt{is\_reply}   & Boolean               & Biến cờ xác định trạng thái: \texttt{True} nếu là bình luận con (phản hồi) và \texttt{False} nếu là bình luận gốc. \\ \hline
\end{tabular}
\end{table}
